% =========================================================================
%  Ejercicio 2b - Ensayo INDIVIDUAL sobre el articulo
%  Responsable: Yahir León Bautista
%
%  NOTA: Este ensayo es INDIVIDUAL, uno por cada integrante del equipo.
% =========================================================================

\newpage
\subsection*{Ejercicio 2b. Ensayo individual}

\begin{flushright}
\textbf{Autor: León Bautista Yahir}
\end{flushright}

\subsubsection*{Introducción}

El artículo analiza la importancia y el valor que tienen los datos personales 
dentro de la economía digital. La autora explica que, con el crecimiento del Big Data, las redes sociales, 
los dispositivos móviles y el Internet de las Cosas (IoT), las personas generamos constantemente información 
que puede ser recopilada, almacenada, procesada y utilizada por empresas para obtener beneficios económicos.
Considero que este tema es especialmente relevante porque muchas veces utilizamos aplicaciones, redes sociales 
y dispositivos conectados sin pensar en la cantidad de información que estamos entregando. Mi postura es que los 
datos personales deben considerarse un recurso de gran valor y que, por ello, las personas debemos tener mayor 
conocimiento y control sobre la información que generamos y compartimos.

\subsubsection*{Desarrollo}

La autora sostiene que nuestros datos personales se han convertido en un recurso de gran valor dentro de la economía digital. 
Cada vez que utilizamos redes sociales, teléfonos, aplicaciones o dispositivos conectados generamos información que puede 
ser almacenada y analizada por las empresas para obtener beneficios, mejorar servicios y crear publicidad personalizada.
El Big Data y el Internet de las Cosas aumentan todavía más esta situación, ya que permiten recopilar enormes cantidades
 de información de manera constante. Sin embargo, la autora señala que muchas personas desconocen qué datos comparten, 
 quién los utiliza y para qué, lo que genera problemas relacionados con la privacidad y el control de la información.
El objetivo principal del artículo es crear conciencia sobre el valor de nuestros datos y defender que los usuarios 
tengan mayor control sobre ellos. La autora propone tecnologías centradas en las personas, como los Personal Data Stores, 
que permitirían administrar y decidir cómo utilizar la información personal.
Esto se relaciona directamente con Fundamentos de Bases de Datos, ya que las bases de datos son fundamentales para almacenar, 
organizar y procesar la información que las empresas recopilan. Como estudiante de Ciencias de la Computación, 
considero importante no solo saber manejar estos datos técnicamente, sino también hacerlo de manera responsable, respetando 
la privacidad y los derechos de los usuarios.
La temática central es el valor de los datos personales en la era del Big Data y el IoT. Como temas secundarios aparecen 
la privacidad, la protección de datos, la publicidad personalizada, el consentimiento y la regulación. Todos ellos son 
relevantes para la práctica profesional, debido a que los sistemas informáticos actuales dependen cada vez más de la 
recopilación y análisis de información.

\subsubsection*{Conclusión}

Después de analizar la lectura, reafirmo mi postura de que nuestros datos personales tienen un valor mucho mayor del 
que normalmente reconocemos. La principal aportación del artículo es mostrar que cada actividad digital puede generar 
información que posteriormente se convierte en un recurso económico, y que por esta razón no debemos considerar nuestros 
datos como algo sin importancia.
Relacionar esta lectura con Fundamentos de Bases de Datos permite comprender que la gestión de información tiene una 
dimensión tanto técnica como ética. Aprendí que almacenar y procesar datos implica también considerar quién tiene acceso 
a ellos, cómo se utilizan y qué derechos tiene la persona a la que pertenecen. 
Finalmente, el crecimiento del Big Data y del Internet de las Cosas plantea preguntas que seguirán siendo relevantes: 
¿quién debería ser el propietario de los datos?, ¿cuánto control debería tener una persona sobre la información que genera?, 
¿qué límites deberían existir para su comercialización? Estas preguntas muestran que el avance tecnológico continuará exigiendo 
un equilibrio entre innovación, aprovechamiento de la información y protección de la privacidad.


% Extension: entre 1 y 2 cuartillas por integrante.
%
% Estructura obligatoria (en parrafos, sin listas):
% - INTRODUCCION: presenta brevemente el tema, por que es importante, y
%   plantea tu postura o idea principal.
% - DESARROLLO: resume los argumentos del autor, relaciona con otros
%   conocimientos/lecturas/ejemplos, expone tus propios argumentos (a favor
%   o en contra) con razones, indica el objetivo que planteo el autor y como
%   se relaciona con la materia, y la tematica central y temas laterales y su
%   relacion con tu practica profesional.
% - CONCLUSION: reafirma tu postura, que aprendiste o la aportacion mas
%   valiosa, y si el tema abre preguntas o retos futuros.
% - FORMATO: parrafos completos, conectores logicos, ortografia y gramatica
%   cuidadas.

% Referencia del articulo (APA) - completar los datos:
% \cite{datavalue}